\chapter{개입으로 메커니즘을 시험하기}
\label{chap:phase6}

\section{결론부터: 위치 찾기와 메커니즘 설명은 다르다}

어떤 activation을 바꾸었을 때 출력이 달라지면 그 activation은 해당 개입과 metric 아래에서 인과적으로 관련된다.
그러나 그것만으로 그 activation이 어떤 정보를 계산했는지, 어느 upstream source에서 왔는지, 어떤 downstream path를 통해 사용되었는지는 알 수 없다.
localization은 메커니즘 연구의 시작점이지 종료점이 아니다.

재현 가능한 개입에는 base input부터 aggregation까지 여섯 요소가 필요하다.
이 장에서는 이를 다음 명세로 기록한다.
\begin{equation}
  \mathcal{I}=(\text{base input},\text{source input},\text{hook},
  \text{replacement},\text{metric},\text{aggregation}).
  \label{eq:intervention-spec}
\end{equation}
``layer 10을 patch했다''는 문장만으로는 residual input인지 output인지, 어느 position과 feature인지, 무엇으로 교체했는지가 빠져 있어 재현할 수 없다.

\section{관찰, 상관과 개입}

관찰 분포와 개입 분포는 서로 다른 질문에 답한다.
계산그래프에서 activation $H$와 출력 $Y$가 함께 변할 때, 관찰 분포 $p(Y\mid H=h)$는 upstream 입력이 두 변수에 함께 영향을 준 결과를 포함한다.
개입은 그래프의 해당 값을 외부에서 지정하고 나머지 downstream 계산을 다시 실행하므로 개념적으로
\begin{equation}
  p(Y\mid \operatorname{do}(H\leftarrow h'))
\end{equation}
를 묻는 셈이다.

신경망의 순전파 그래프는 관측 변수 사이의 일반적인 사회과학적 인과 문제보다 직접 조작하기 쉽다.
다만 개입한 activation이 훈련 중 나타나지 않은 값이면 downstream module이 distribution 밖에서 작동한다.
따라서 인과 효과는 ``원래 모델이 자연스럽게 그 상태를 만드는 과정''과 ``강제로 그 상태를 넣었을 때의 반응''을 구분하여 해석해야 한다.

\begin{factbox}
activation intervention은 parameter를 다시 학습하지 않고 한 번의 실행에서 중간값을 교체한다.
weight editing은 이후 모든 입력에 쓰이는 함수를 바꾼다.
두 실험은 서로 다른 인과 질문에 답한다.
\end{factbox}

\section{ablation}

ablation은 구성요소의 정상 출력을 기준값으로 바꾸거나 특정 edge를 끊고 target metric을 측정한다.
component $c$의 정상 실행 metric을 $m_{\mathrm{full}}$, ablated 실행을 $m_{-c}$라고 하면 단순 효과는
\begin{equation}
  \Delta_c=m_{\mathrm{full}}-m_{-c}
\end{equation}
로 쓸 수 있다.
metric의 방향에 따라 부호 정의를 바꿔야 한다.
\index{ablation@절제(ablation)}

중요한 것은 ``제거''가 수학적으로 유일하지 않다는 점이다.

\begin{description}[style=nextline,leftmargin=2em]
  \item[zero ablation]
  activation을 0으로 바꾼다.
  0이 해당 hook에서 자연스러운 baseline인지 확인해야 한다.

  \item[mean ablation]
  reference dataset의 평균 activation으로 바꾼다.
  평균을 계산한 조건에 따라 다른 정보를 보존할 수 있다.

  \item[resample ablation]
  다른 표본의 activation을 넣는다.
  marginal distribution은 더 잘 보존할 수 있지만 source 표본의 정보가 새로 들어간다.

  \item[noise ablation]
  선택한 noise 분포를 더하거나 대체한다.
  scale과 covariance가 결과를 결정한다.
\end{description}

head output을 0으로 만드는 것과 attention score를 $-\infty$로 바꾸어 특정 source를 막는 것은 다른 개입이다.
전자는 head의 모든 source contribution을 없애고, 후자는 softmax가 나머지 source에 weight를 재분배한다.

\section{activation patching과 causal tracing}

activation patching은 서로 다른 두 입력의 실행을 결합한다.
clean input $x_c$에서는 원하는 행동이 나타나고, corrupted input $x_r$에서는 그 행동이 약해진다고 하자.
먼저 두 실행의 activation을 cache한다.
그다음 $x_r$을 실행하면서 선택한 hook $u$의 값을 clean activation $h_u(x_c)$로 교체한다.
\index{activation patching@활성값 패칭(activation patching)}

patch recovery는 corrupted 실행이 clean activation의 이식으로 얼마나 회복되는지를 정규화한다.
metric $m$에 대해
\begin{equation}
  R_u=\frac{m(x_r;h_u\leftarrow h_u(x_c))-m(x_r)}
  {m(x_c)-m(x_r)}
  \label{eq:patch-recovery}
\end{equation}
로 정의할 수 있다.
분모가 작으면 불안정하며, $R_u$는 0과 1 사이에 반드시 머물지 않는다.
확률, logit, logit difference와 loss 가운데 무엇을 $m$으로 쓰는지에 따라 결과가 달라질 수 있다 \citep{zhang2024activation,heimersheim2024patching}.

ROME 연구의 causal tracing은 factual prompt의 subject embedding에 noise를 넣어 prediction을 손상시킨 뒤, 특정 layer와 position의 hidden activation을 clean 값으로 복원하여 사실 예측이 얼마나 회복되는지 측정했다.
이 실험은 GPT 계열 모델에서 subject token을 처리하는 middle-layer feed-forward computation의 역할을 지지했다 \citep{meng2022rome}.

\begin{limitbox}
patch가 큰 회복을 만들었다는 사실은 해당 hook이 clean과 corrupted 조건의 차이를 운반하는 충분한 매개 지점임을 시사한다.
정보가 그 layer에서 처음 계산되었다거나 오직 그 위치에 저장되었다는 사실까지 증명하지는 않는다.
residual stream은 upstream 차이를 누적해서 운반할 수 있다.
\end{limitbox}

\section{path patching과 edge intervention}

node patching은 한 activation으로 들어오는 모든 upstream path를 함께 바꾼다.
path 또는 edge patching은 sender $s$의 clean/corrupted 차이가 특정 receiver $r$에 들어가는 경로만 바꾸고, 다른 입력은 원 조건으로 유지하려고 한다.
예를 들어 head output이 뒤 head의 key input으로 가는 edge만 교체하면 value와 query 경로를 분리해 시험할 수 있다.
\index{path patching@경로 패칭(path patching)}

Transformer 구현에서는 residual stream에 여러 sender output이 더해진 뒤 receiver가 이를 읽는다.
따라서 edge를 직접 저장한 tensor가 없을 수 있다.
연구자는 receiver input을 sender별 기여로 재구성하거나 여러 차례 forward pass로 원하는 경로만 patch해야 한다.
어떤 tensor decomposition을 사용했는지가 곧 edge의 정의다.

path patching의 결과도 상호작용에 민감하다.
두 경로가 서로 대체 가능하면 각각을 따로 끊을 때 효과가 작지만 함께 끊을 때 효과가 클 수 있다.
반대로 한 경로의 효과가 다른 경로가 존재할 때만 나타날 수 있다.
따라서 single-edge effect를 단순히 합하여 전체 효과를 얻을 수 있다고 가정하면 안 된다.

\section{interchange intervention과 causal abstraction}

interchange intervention은 source example에서 얻은 내부 상태를 base example의 계산에 이식하여, 고수준 변수의 값을 바꾼 것과 같은 반사실 출력이 나오는지 검사한다.
고수준 causal model의 변수 $Z$와 신경망 activation 집합 $H$ 사이에 alignment를 제안했다고 하자.
source $s$의 $Z$ 값을 base $b$에 넣은 고수준 출력과, $H(s)$를 $b$의 신경망에 넣은 저수준 출력이 일치하는지를 여러 pair에서 검사한다.
\index{causal abstraction@인과 추상화(causal abstraction)}

Geiger 등은 이 생각을 causal abstraction의 형식으로 정리하고, MQNLI를 위해 구성된 자연논리 causal model의 일부 구조가 BERT 기반 모델에서 실현되는지를 interchange intervention으로 시험했다 \citep{geiger2021causal}.
이 방법의 강점은 ``이 neuron이 negation을 뜻한다'' 같은 이름 붙이기를 넘어, 제안한 고수준 변수와 동일한 반사실 관계를 만드는지 검증한다는 데 있다.

causal abstraction은 연구자가 먼저 고수준 causal model, 변수와 alignment 후보를 정해야 한다는 한계가 있다.
시험을 통과하면 그 abstraction이 유효하다는 증거지만 유일한 abstraction이라는 뜻은 아니다.
자연어 전반에 대한 올바른 고수준 변수를 자동으로 얻는 문제도 해결되지 않는다.

\section{attribution과 gradient 근사}

first-order attribution은 작은 activation 변화의 metric 효과를 gradient 내적으로 근사한다.
미분 가능한 metric $m$과 activation $\vect{h}$에 대해 변화 $\Delta\vect{h}$를 주면
\begin{equation}
  m(\vect{h}+\Delta\vect{h})-m(\vect{h})
  \approx \nabla_{\vect{h}}m^{\mathsf T}\Delta\vect{h}
  \label{eq:first-order-attribution}
\end{equation}
다.
attribution patching은 clean과 corrupted activation 차이를 $\Delta\vect{h}$로 두어 실제 patch를 반복하지 않고 많은 node나 edge의 효과를 근사한다.

Syed 등은 두 번의 forward pass와 한 번의 backward pass를 사용하는 이 선형 근사로 edge importance를 계산하고, 조사한 benchmark에서 기존 자동 회로 발견법보다 평균적인 circuit-recovery AUC가 높았다고 보고했다 \citep{syed2023eap}.
계산량 절감은 중요하지만 근사가 큰 intervention, activation threshold와 상호작용을 정확히 포착한다는 보장은 없다.

gradient가 0에 가깝다고 구성요소가 불필요하다고 단정할 수도 없다.
saturation, canceling path 또는 first-order에서 보이지 않는 곡률이 있을 수 있다.
반대로 큰 gradient는 매우 작은 국소 변화에 민감하다는 뜻이며 자연스러운 크기의 patch 효과와 같지 않을 수 있다.

\section{자동 회로 발견}

자동 회로 발견은 모든 edge를 사람이 직접 검사하는 비용을 줄이려는 방법이다.
ACDC는 정한 dataset과 metric 아래에서 edge를 반복적으로 제거하고 성능 변화를 기준으로 graph를 prune한다.
GPT-2 small의 greater-than 과제에서 이전 수작업 연구가 찾은 5개 component type을 모두 다시 찾고, 약 32,000개 edge 가운데 68개를 선택한 결과가 보고되었다 \citep{conmy2023acdc}.

자동화가 해결하는 것은 주어진 node/edge space에서 후보 subgraph를 찾는 검색 문제다.
다음 항목은 여전히 사람이 정의하거나 검증해야 한다.

\begin{itemize}
  \item 어떤 behavior와 input distribution을 설명할 것인가?
  \item node와 edge를 어느 granularity로 나눌 것인가?
  \item 어떤 ablation baseline과 metric을 사용할 것인가?
  \item 선택된 구성요소를 어떤 algorithmic role로 설명할 것인가?
  \item 새로운 반사실 입력에서 회로가 예측을 맞히는가?
\end{itemize}

2026년의 CircuitLasso는 high-dimensional SAE feature 회로를 sparse linear regression으로 학습하여 intervention 기반 방법보다 적은 비용으로 확장하려는 접근을 제시했다.
저자들은 benchmark에서 기존 intervention 방법과 비슷한 structural accuracy를 보고했지만, 이는 2026년 6월 공개된 연구의 평가 범위에 한정된다 \citep{yin2026circuitlasso}.

\section{faithfulness, completeness, minimality}

회로 $C$와 전체 모델 $M$을 비교할 때 세 기준을 다음처럼 구분할 수 있다.
\index{faithfulness@충실성(faithfulness)}

\begin{description}[style=nextline,leftmargin=2em]
  \item[faithfulness]
  회로 밖을 선택한 baseline으로 ablate했을 때 $C$가 전체 모델의 target behavior를 얼마나 보존하는가?

  \item[completeness]
  중요한 계산이 회로 밖에 남아 있지 않은가?
  회로의 complement 또는 추가 구성요소가 설명되지 않은 성능을 만들지 않는지 검사한다.

  \item[minimality]
  회로 안의 각 요소가 필요하거나, 제거하면 정한 기준 이상으로 설명력이 감소하는가?
\end{description}

이 개념에는 ablation 이후 모델을 어떻게 정의하는지가 들어간다.
Miller 등은 circuit faithfulness score가 겉보기에 작은 ablation 방법 변화에도 크게 달라질 수 있음을 보고했다 \citep{miller2024faithfulness}.
따라서 숫자 하나를 회로의 내재적 속성처럼 보고하면 안 된다.
최소한 baseline distribution, patch 단위, metric, normalization과 회로 크기 곡선을 함께 제시해야 한다.

\begin{studybox}[회로 평가의 권장 보고]
회로 평가는 보존 성능, 제거 효과, 크기 일치 대조군과 baseline 민감도를 함께 보고해야 한다.
전체 모델, 회로만 남긴 모델, 회로를 제거한 모델과 무작위 크기 일치 회로를 같은 표에 둔다.
zero, mean과 resample ablation 가운데 가능한 둘 이상을 비교한다.
개별 sample 분포와 bootstrap confidence interval을 보고하고, 회로 크기에 따른 metric curve를 제시한다.
\end{studybox}

\section{self-repair와 방법 민감도}

구성요소를 ablate하면 downstream layer는 원래 실행에서 보이지 않던 보상 update를 만들 수 있다.
Hydra effect 연구는 한 attention layer를 제거했을 때 다른 layer가 일부 효과를 보상하는 사례와 late MLP가 최대가능도 token을 낮추는 counterbalancing pattern을 보고했다 \citep{mcgrath2023hydra}.
이때 ablation 후 최종 logit 변화는 제거한 구성요소의 원래 direct contribution보다 작을 수 있다.
\index{self-repair@자기 복구(self-repair)}

작은 ablation 효과에는 구성요소의 작은 기여 외에도 여러 대안 설명이 있다.
구체적으로 다음 원인 가운데 하나 이상이 결과를 만들 수 있다.

\begin{itemize}
  \item 원래 구성요소의 기여가 실제로 작다.
  \item 병렬 중복 경로가 같은 기능을 한다.
  \item ablation 때문에 downstream 보상이 새로 발생한다.
  \item metric이 해당 기여를 잘 드러내지 못한다.
  \item baseline이 원 activation과 우연히 비슷한 기능을 보존한다.
\end{itemize}

activation patching의 corruption 방식과 metric만 바꾸어도 localization map이 달라질 수 있다는 체계적 비교가 보고되었다 \citep{zhang2024activation}.
그러므로 한 heatmap을 메커니즘 자체처럼 해석하지 말고, 합리적인 대안 설정에서 결론이 얼마나 유지되는지 sensitivity analysis를 수행해야 한다.

\section{Phase 6 학습 점검}

\begin{studybox}
\begin{enumerate}
  \item clean/corrupted prompt pair를 만들고 두 조건의 차이가 목표 변수 하나에 가깝도록 검토한다.
  \item residual pre, attention output, MLP output을 position과 layer별로 patch하여 logit-difference recovery map을 만든다.
  \item 같은 실험을 probability와 loss metric으로 반복하고 localization 순위를 비교한다.
  \item zero, dataset mean과 resample ablation을 비교하여 baseline 민감도를 기록한다.
  \item node patching으로 찾은 후보를 path patching으로 좁히고, sender와 receiver의 구체적인 알고리즘 역할을 반사실 prompt로 시험한다.
  \item 발견한 회로의 faithfulness, completeness와 minimality를 서로 다른 표에 보고한다.
\end{enumerate}
\end{studybox}

\begin{factbox}
Phase 6의 종료 조건은 intervention 코드를 실행할 수 있다는 것이 아니다.
개입이 답하는 인과 질문, 개입이 만드는 distribution shift, metric과 baseline 의존성, 그리고 결과가 localization인지 mechanism인지 정확히 구분할 수 있어야 한다.
\end{factbox}
